% !TeX program = pdflatex
\documentclass[11pt,a4paper]{amsart}

% Shared preamble for course notes and exercise sets.

\usepackage[T1]{fontenc}
\usepackage{lmodern}
\usepackage[english]{babel}

\usepackage[a4paper,margin=30mm]{geometry}
\usepackage{microtype}
\usepackage{graphicx}
\usepackage{enumitem}

\usepackage{amsmath}
\usepackage{amssymb}
\usepackage{amsthm}
\usepackage{mathtools}
\usepackage{aliascnt}

\usepackage[hidelinks,pdfusetitle]{hyperref}

\numberwithin{equation}{section}
\setlist{itemsep=0.25em,topsep=0.5em}

\theoremstyle{plain}
\newtheorem{theorem}{Theorem}[section]
\newaliascnt{lemma}{theorem}
\newtheorem{lemma}[lemma]{Lemma}
\aliascntresetthe{lemma}
\newaliascnt{proposition}{theorem}
\newtheorem{proposition}[proposition]{Proposition}
\aliascntresetthe{proposition}
\newaliascnt{corollary}{theorem}
\newtheorem{corollary}[corollary]{Corollary}
\aliascntresetthe{corollary}

\theoremstyle{definition}
\newaliascnt{definition}{theorem}
\newtheorem{definition}[definition]{Definition}
\aliascntresetthe{definition}
\newaliascnt{example}{theorem}
\newtheorem{example}[example]{Example}
\aliascntresetthe{example}
\newaliascnt{problem}{theorem}
\newtheorem{problem}[problem]{Problem}
\aliascntresetthe{problem}

\theoremstyle{remark}
\newaliascnt{remark}{theorem}
\newtheorem{remark}[remark]{Remark}
\aliascntresetthe{remark}

\usepackage[nameinlink,noabbrev]{cleveref}

\crefname{theorem}{theorem}{theorems}
\crefname{lemma}{lemma}{lemmas}
\crefname{proposition}{proposition}{propositions}
\crefname{corollary}{corollary}{corollaries}
\crefname{definition}{definition}{definitions}
\crefname{example}{example}{examples}
\crefname{problem}{problem}{problems}
\crefname{remark}{remark}{remarks}

\DeclareMathOperator{\im}{im}
\DeclareMathOperator{\rank}{rank}
\DeclarePairedDelimiter{\abs}{\lvert}{\rvert}
\DeclarePairedDelimiter{\norm}{\lVert}{\rVert}
\DeclarePairedDelimiter{\set}{\lbrace}{\rbrace}


\usepackage{tikz}

\title{Introduction to University Mathematics}
\author{Joona Piirainen}
\date{Autumn 2026}

\begin{document}

\maketitle

\begin{abstract}
  Notes for the University of Helsinki course MAT11001, taught in Finnish
  from 31.08.2026--17.10.2026.
\end{abstract}

\tableofcontents

\section{Course information}

The official course information is available on the
\href{https://studies.helsinki.fi/kurssit/toteutus/hy-opt-cur-2627-8d4269d2-7e32-4dce-bcdb-b7fa10fecac6}{University
of Helsinki course page}.

\section{Lecture notes}

\subsection{31.08.2026}

The first lecture was mostly spent going over the course structure and other
general admin topics. We did go over some preliminary definitions and examples,
which will follow.

\subsubsection{Notational conventions}

\begin{definition}
  An integer \(p > 1\) is \emph{prime} if its only positive divisors are
  \(1\) and \(p\).
\end{definition}

\begin{definition}
  The following basic sets of numbers will be used throughout:
  \begin{align*}
    \mathbb{N} &= \{0, 1, 2, \ldots\} \\
    \mathbb{Z} &= \{\ldots, -2, -1, 0, 1, 2, \ldots\} \\
    \mathbb{Z}^{+} &= \mathbb{N} \setminus \{0\} \\
    \mathbb{Q} &= \left\{\frac{p}{q} : p,q \in \mathbb{Z},\ q \neq 0\right\} \\
    \mathbb{R} &= \text{the set of real numbers} \\
    \mathbb{C} &= \left\{a+bi : a,b \in \mathbb{R}\right\}, \quad i^2=-1 \\
    \mathbb{P} &= \left\{p \in \mathbb{Z} : p \text{ is prime}\right\}
  \end{align*}
\end{definition}

\subsubsection{Principle of mathematical induction}

Here we state the fundamental principle of mathematical induction without proof,
and use it freely thereafter.

\begin{theorem}[Principle of mathematical induction]\label{thm:induction}
  Let \(n_0 \in \mathbb{N}\), and let \(P(n)\) be a statement for each
  \(n \in \mathbb{N}\) with \(n \geq n_0\). Suppose that
  \begin{enumerate}
    \item \(P(n_0)\) is true, and
    \item for every \(n \geq n_0\), \(P(n)\) implies \(P(n+1)\).
  \end{enumerate}
  Then \(P(n)\) is true for every \(n \geq n_0\).
\end{theorem}

\begin{proof}
  Omitted.
\end{proof}

\begin{example}
  We prove by induction that
  \[
    \sum_{k=1}^{n} k = \frac{n(n+1)}{2}
  \]
  for every integer \(n \geq 1\).

  For \(n=1\), the identity holds because
  \[
    \sum_{k=1}^{1} k = 1 = \frac{1(1+1)}{2}.
  \]
  Now suppose that the identity holds for some integer \(n \geq 1\). Then
  \begin{align*}
    \sum_{k=1}^{n+1} k
    &= \sum_{k=1}^{n} k + (n+1) \\
    &= \frac{n(n+1)}{2} + (n+1) \\
    &= \frac{(n+1)(n+2)}{2}.
  \end{align*}
  Thus the identity also holds for \(n+1\), and hence for every integer
  \(n \geq 1\). \qed
\end{example}

\subsection{03.09.2026}

\subsubsection{Divisibility of integers}

We start by defining the precise meaning of \emph{even} and \emph{odd} numbers.

\begin{definition}
  A number \(m \in \ZZ\) is \emph{even} if there exists \(k \in \ZZ\) such that
  \[
    m = 2k.
  \]
  On the other hand, \(m\) is \emph{odd} if there exists a \(k \in \ZZ\) such
  that
  \[
    m = 2k + 1.
  \]
\end{definition}

\begin{definition}[Divisibility]
  A number \(a \in \ZZ\) is said to be \emph{divisible} by a number
  \(b \in \ZZ, b \neq 0\), if there exists a \(k \in \ZZ\) such that
  \[
    a = kb.
  \]
  In this case we also say that \(b\) \emph{divides} \(a\), and write
  \[
    b \mid a.
  \]
  If \(b\) does \emph{not} divide \(a\), we write
  \[
    b \nmid a.
  \]
\end{definition}

\begin{definition}
  For \(n \in \ZZ \setminus \{0\}\), we denote the set of integers that are
  divisible by \(n\) by
  \[
    n\ZZ = \left\{z \in \ZZ : z = nk, \quad k \in \ZZ\right\}.
  \]
\end{definition}

\subsubsection{Strong induction}

In some proofs it is useful to assume all preceding cases rather than only the
immediately preceding case. This gives an equivalent form of the principle of
induction, called \emph{strong induction}.

\begin{theorem}[Strong induction]\label{thm:strong-induction}
  Let \(n_0 \in \NN\), and let \(P(n)\) be a statement for each
  \(n \in \NN\) with \(n \geq n_0\). Suppose that
  \begin{enumerate}
    \item \(P(n_0)\) is true, and
    \item for every \(n \geq n_0\), if \(P(k)\) is true for every \(k\) with
      \(n_0 \leq k \leq n\), then \(P(n + 1)\) is true.
  \end{enumerate}
  Then \(P(n)\) is true for every \(n \geq n_0\).
\end{theorem}

\begin{proof}
  Omitted.
\end{proof}

As an example where strong induction is natural, let us prove that all natural
numbers \(n \geq 2\) can be written as a product of \emph{prime} numbers.

\begin{theorem}
  Every natural number \(n \geq 2\) can be written as a product of
  prime numbers.
\end{theorem}

\begin{proof}
  When \(n = 2\), we have \(n \in \PP\), so the result holds trivially.
  Now suppose \(n \geq 2\) and the statement holds for all
  \(2 \leq k \leq n\). If \(n + 1 \in \PP\), the result holds trivially, so
  suppose \(n + 1 \notin \PP\). Then we have
  \[
    n + 1 = ab,
  \]
  for some \(2 \leq a,b \leq n\). By the induction hypothesis, \(a\) and
  \(b\) can be written as products of primes, so
  \[
    \begin{aligned}
      a &= p_1p_2\cdots p_i, \quad p_1,\ldots,p_i \in \PP, \\
      b &= q_1q_2\cdots q_j, \quad q_1,\ldots,q_j \in \PP.
    \end{aligned}
  \]
  Thus we have
  \[
    n + 1 = ab = p_1p_2\cdots p_iq_1q_2\cdots q_j.
  \]
  The result follows from Theorem \ref{thm:strong-induction}.
\end{proof}

\subsubsection{Sequences}

\begin{definition}
  A sequence of real numbers is a function \(z\colon \NN \rightarrow \RR\),
  whose terms are denoted by \(z_n = z(n)\). We write the sequence as
  \[
    (z_n)_{n \in \NN} = (z_0,z_1,z_2,\ldots).
  \]
\end{definition}

\begin{definition}
  A sequence of real numbers may also be specified recursively. First, for
  some \(k \in \ZZ^+\), we define the initial terms
  \(z_0,\ldots,z_{k-1}\). For each \(n \geq k\), we then define \(z_n\) in
  terms of earlier terms \(z_0,\ldots,z_{n-1}\).
\end{definition}

\subsection{07.09.2026}

This lecture was mostly spent going through basic mathematical notation.

\subsubsection{Logical connectives}

\begin{definition}
  Let \(p\) and \(q\) be propositions. We define the following logical
  connectives:
  \begin{enumerate}
    \item The \emph{conjunction} \(p \land q\), read ``\(p\) and \(q\)'', is
      true exactly when both \(p\) and \(q\) are true.
    \item The \emph{disjunction} \(p \lor q\), read ``\(p\) or \(q\)'', is
      true exactly when at least one of \(p\) and \(q\) is true.
    \item The \emph{implication} \(p \implies q\), read ``\(p\) implies
      \(q\)'', is false exactly when \(p\) is true and \(q\) is false.
    \item The \emph{equivalence} \(p \iff q\), read ``\(p\) if and only if
      \(q\)'', is true exactly when \(p\) and \(q\) have the same truth value.
    \item The \emph{negation} \(\neg p\), read ``not \(p\)'', is true exactly
      when \(p\) is false.
  \end{enumerate}
\end{definition}

\begin{definition}
  Let \(p\) and \(q\) be propositions.
  \begin{enumerate}
    \item The \emph{negated implication} \(p \nRightarrow q\), read ``\(p\)
      does not imply \(q\)'', is true exactly when \(p\) is true and \(q\) is
      false.
    \item The \emph{negated equivalence} \(p \nLeftrightarrow q\), read
      ``\(p\) is not equivalent to \(q\)'', is true exactly when \(p\) and
      \(q\) have different truth values.
  \end{enumerate}
\end{definition}

\begin{definition}
  Let \(X\) be a set. A \emph{predicate} \(P(x)\) on \(X\) is a statement
  containing a variable \(x\) whose truth value depends on the chosen value of
  \(x \in X\). For each \(a \in X\), substituting \(a\) for \(x\) gives the
  proposition \(P(a)\), which is either true or false. For example,
  \(P(x) \coloneqq x > 0\) is a predicate on \(\RR\); \(P(2)\) is true, whereas
  \(P(-1)\) is false.
\end{definition}

\begin{definition}
  Let \(P(x)\) be a predicate on a set \(X\).
  \begin{enumerate}
    \item The \emph{universal quantification}
      \(\forall x \in X \colon P(x)\), read ``for every \(x \in X\),
      \(P(x)\)'', is true exactly when \(P(a)\) is true for every \(a \in X\).
    \item The \emph{existential quantification}
      \(\exists x \in X \colon P(x)\), read ``there exists an \(x \in X\) such
      that \(P(x)\)'', is true exactly when \(P(a)\) is true for at least one
      \(a \in X\).
  \end{enumerate}
\end{definition}

\begin{definition}
  The \emph{negated existential quantification}
  \(\nexists x \in X \colon P(x)\), read ``there does not exist an \(x \in X\)
  such that \(P(x)\)'', is true exactly when \(P(a)\) is false for every
  \(a \in X\). Equivalently,
  \[
    \nexists x \in X \colon P(x)
    \quad \iff \quad
    \forall x \in X \colon \neg P(x).
  \]
\end{definition}

\subsubsection{Basic set theory}

\begin{definition}
  A \emph{set} is a collection of distinct objects, called its
  \emph{elements}. If an object \(x\) is an element of a set \(A\), we write
  \(x \in A\); otherwise, we write \(x \notin A\). The order in which the
  elements are listed, and any repetition among them, do not affect the set.
\end{definition}

\begin{definition}
  Two sets \(A\) and \(B\) are \emph{equal}, written \(A = B\), if they have
  exactly the same elements. Equivalently,
  \[
    A = B \iff \forall x \colon (x \in A \iff x \in B).
  \]
\end{definition}

\begin{definition}
  The \emph{empty set} is the set containing no elements. It is denoted by
  \(\emptyset\).
\end{definition}

\begin{definition}
  Let \(A\) and \(B\) be are sets. \(A\) and \(B\) are said to be
  \emph{disjoint}, if \(A \cap B = \emptyset\).
\end{definition}

\begin{definition}
  A \emph{singleton} set is a set consisting of exactly \textbf{one} element.
\end{definition}

\begin{definition}
  Let \(P(x)\) be a predicate on a set \(X\). The \emph{set-builder notation}
  \[
    A \coloneqq \set{x \in X \colon P(x)}
  \]
  defines \(A\) to be the set of precisely those elements \(x \in X\) for which
  \(P(x)\) is true. Equivalently, for every \(a \in X\),
  \[
    a \in A \iff P(a).
  \]
\end{definition}

\begin{definition}
  A set \(A\) is a \emph{subset} of a set \(B\) if every element of \(A\) is
  also an element of \(B\). We then write \(A \subseteq B\). If \(A\) is not a
  subset of \(B\), we write \(A \nsubseteq B\); this means that some element of
  \(A\) is not an element of \(B\). If \(A \subseteq B\) and \(A \neq B\), then
  \(A\) is a \emph{proper subset} of \(B\), written \(A \subset B\).
  Two sets are equal exactly when both containments hold:
  \[
    A = B \iff A \subseteq B \land B \subseteq A.
  \]
\end{definition}

\begin{definition}
  Let \(A\) and \(B\) be subsets of a set \(X\).
  \begin{enumerate}
    \item The \emph{union} of \(A\) and \(B\) is
      \[
        A \cup B \coloneqq
        \set{x \in X \colon x \in A \lor x \in B}.
      \]
    \item The \emph{intersection} of \(A\) and \(B\) is
      \[
        A \cap B \coloneqq
        \set{x \in X \colon x \in A \land x \in B}.
      \]
    \item The \emph{difference} of \(A\) and \(B\) is
      \[
        A \setminus B \coloneqq
        \set{x \in X \colon x \in A \land x \notin B}.
      \]
    \item The \emph{complement} of \(A\) in \(X\) is
      \[
        \complement{A} \coloneqq X \setminus A
        = \set{x \in X \colon x \notin A}.
      \]
    \item The \emph{symmetric difference} of \(A\) and \(B\) is
      \[
        A \mathbin{\Delta} B \coloneqq
        (A \setminus B) \cup (B \setminus A).
      \]
  \end{enumerate}
\end{definition}

\begin{center}
  \begin{tikzpicture}[scale=0.9, every node/.style={font=\small}]
    % Union
    \begin{scope}
      \fill[blue!25] (1.4,1.2) circle (0.85);
      \fill[blue!25] (2.2,1.2) circle (0.85);
      \draw (0,0) rectangle (3.6,2.4);
      \draw (1.4,1.2) circle (0.85);
      \draw (2.2,1.2) circle (0.85);
      \node at (0.2,2.18) {$X$};
      \node at (1.05,1.2) {$A$};
      \node at (2.55,1.2) {$B$};
      \node[below] at (1.8,0) {$A \cup B$};
    \end{scope}

    % Intersection
    \begin{scope}[xshift=5cm]
      \begin{scope}
        \clip (1.4,1.2) circle (0.85);
        \fill[blue!25] (2.2,1.2) circle (0.85);
      \end{scope}
      \draw (0,0) rectangle (3.6,2.4);
      \draw (1.4,1.2) circle (0.85);
      \draw (2.2,1.2) circle (0.85);
      \node at (0.2,2.18) {$X$};
      \node at (1.05,1.2) {$A$};
      \node at (2.55,1.2) {$B$};
      \node[below] at (1.8,0) {$A \cap B$};
    \end{scope}

    % Difference
    \begin{scope}[yshift=-3.4cm]
      \begin{scope}
        \clip (1.4,1.2) circle (0.85);
        \fill[blue!25] (0,0) rectangle (3.6,2.4);
        \fill[white] (2.2,1.2) circle (0.85);
      \end{scope}
      \draw (0,0) rectangle (3.6,2.4);
      \draw (1.4,1.2) circle (0.85);
      \draw (2.2,1.2) circle (0.85);
      \node at (0.2,2.18) {$X$};
      \node at (1.05,1.2) {$A$};
      \node at (2.55,1.2) {$B$};
      \node[below] at (1.8,0) {$A \setminus B$};
    \end{scope}

    % Complement
    \begin{scope}[xshift=5cm,yshift=-3.4cm]
      \fill[blue!25] (0,0) rectangle (3.6,2.4);
      \fill[white] (1.8,1.2) circle (0.85);
      \draw (0,0) rectangle (3.6,2.4);
      \draw (1.8,1.2) circle (0.85);
      \node at (0.2,2.18) {$X$};
      \node at (1.8,1.2) {$A$};
      \node[below] at (1.8,0) {$\complement{A}$};
    \end{scope}

    % Symmetric difference
    \begin{scope}[xshift=2.5cm,yshift=-6.8cm]
      \fill[blue!25] (1.4,1.2) circle (0.85);
      \fill[blue!25] (2.2,1.2) circle (0.85);
      \begin{scope}
        \clip (1.4,1.2) circle (0.85);
        \fill[white] (2.2,1.2) circle (0.85);
      \end{scope}
      \draw (0,0) rectangle (3.6,2.4);
      \draw (1.4,1.2) circle (0.85);
      \draw (2.2,1.2) circle (0.85);
      \node at (0.2,2.18) {$X$};
      \node at (1.05,1.2) {$A$};
      \node at (2.55,1.2) {$B$};
      \node[below] at (1.8,0) {$A \mathbin{\Delta} B$};
    \end{scope}
  \end{tikzpicture}
\end{center}

\subsection{10.09.2026}

\subsubsection{Power sets}

\begin{definition}
  The \emph{power set} of a set \(A\), denoted by \(\powerset{A}\), is the set
  of all subsets of \(A\):
  \[
    \powerset{A} \coloneqq \set{B : B \subseteq A}.
  \]
  In particular, \(\emptyset \in \powerset{A}\) and
  \(A \in \powerset{A}\).
\end{definition}

\begin{definition}
  Let \(A\) and \(B\) be sets. Their \emph{Cartesian product} is the set of all
  ordered pairs whose first component belongs to \(A\) and whose second
  component belongs to \(B\):
  \[
    A \times B \coloneqq \set{(a,b) : a \in A \land b \in B}.
  \]
  More generally, for sets \(A_1,\ldots,A_n\), their Cartesian product is
  \[
    A_1 \times \cdots \times A_n
    \coloneqq
    \set{(a_1,\ldots,a_n) : a_i \in A_i
    \text{ for every } i \in \set{1,\ldots,n}}.
  \]
  If all the factors are the same set \(A\), one commonly writes
  \[
    A^n \coloneqq \underbrace{A \times \cdots \times A}_{n\text{ factors}}.
  \]
\end{definition}

\begin{definition}
  Let \(I\) be a countably infinite \emph{index set}. A family of subsets of a
  set \(X\), indexed by \(I\), is denoted by \((A_i)_{i \in I}\). Its
  \emph{union} and \emph{intersection} are defined by
  \begin{align*}
    \bigcup_{i \in I} A_i
    &\coloneqq \set{x \in X \colon \exists i \in I,\ x \in A_i}, \\
    \bigcap_{i \in I} A_i
    &\coloneqq \set{x \in X \colon \forall i \in I,\ x \in A_i}.
  \end{align*}
  Thus, an element belongs to the union if it belongs to at least one of the
  sets, and it belongs to the intersection if it belongs to every one of them.

  If the family itself is denoted by
  \(\mathcal{A} = \set{A_i : i \in I}\), one also writes
  \(\bigcup \mathcal{A}\) and \(\bigcap \mathcal{A}\).

  A common index set is \(I = \mathbb{N}\). Since \(0 \in \mathbb{N}\) by
  our convention, one writes
  \begin{align*}
    \bigcup_{n=0}^{\infty} A_n
    &= \bigcup_{n \in \mathbb{N}} A_n, &
    \bigcap_{n=0}^{\infty} A_n
    &= \bigcap_{n \in \mathbb{N}} A_n.
  \end{align*}
  For a family indexed by the positive integers, the corresponding notation is
  \(\bigcup_{n=1}^{\infty} A_n\) and
  \(\bigcap_{n=1}^{\infty} A_n\). When the index range is clear from context,
  these are also abbreviated as \(\bigcup_n A_n\) and \(\bigcap_n A_n\).
\end{definition}

\subsubsection{Laws of set algebra}

Let \(A\), \(B\), and \(C\) be subsets of a universal set \(X\). The union
and intersection operations are \emph{commutative}:
\begin{align*}
  A \cup B &= B \cup A, &
  A \cap B &= B \cap A.
\end{align*}
They are also \emph{associative}:
\begin{align*}
  (A \cup B) \cup C &= A \cup (B \cup C), &
  (A \cap B) \cap C &= A \cap (B \cap C).
\end{align*}
The \emph{distributive laws} are
\begin{align*}
  A \cap (B \cup C) &= (A \cap B) \cup (A \cap C), \\
  A \cup (B \cap C) &= (A \cup B) \cap (A \cup C).
\end{align*}
The \emph{De Morgan laws} relate complements to unions and intersections:
\begin{align*}
  \complement{(A \cup B)} &= \complement{A} \cap \complement{B}, \\
  \complement{(A \cap B)} &= \complement{A} \cup \complement{B}.
\end{align*}
Finally, set difference can be expressed using intersection and complement:
\[
  A \setminus B = A \cap \complement{B}.
\]
Equivalently, replacing \(B\) by \(C\) gives
\[
  A \setminus C = A \cap \complement{C}.
\]

\begin{theorem}
  Let \(A\) and \(B\) be subsets of a universal set \(X\). Then
  \[
    A \subseteq B \iff \complement{B} \subseteq \complement{A}.
  \]
\end{theorem}

\begin{proof}
  Suppose first that \(A \subseteq B\), and let
  \(x \in \complement{B}\). Then \(x \notin B\). Since every element of
  \(A\) belongs to \(B\), it follows that \(x \notin A\), and hence
  \(x \in \complement{A}\). Therefore
  \(\complement{B} \subseteq \complement{A}\).

  Conversely, suppose that
  \(\complement{B} \subseteq \complement{A}\), and let \(x \in A\). If
  \(x \notin B\), then \(x \in \complement{B}\), so the assumed inclusion
  would imply that \(x \in \complement{A}\), contradicting \(x \in A\).
  Thus \(x \in B\), and therefore \(A \subseteq B\).
\end{proof}

\subsection{14.09.2026}

\subsubsection{Proof by counterexample}

To disprove a universal claim, it is enough to find a single case in which the
claim fails. Such a case is called a \emph{counterexample}.

\begin{example}
  Consider the claim that
  \[
    (a+b)^2 = a^2+b^2
  \]
  for all real numbers \(a\) and \(b\). Let \(a=b=1\). Then
  \[
    (a+b)^2 = (1+1)^2 = 4,
    \qquad
    a^2+b^2 = 1^2+1^2 = 2.
  \]
  Since \(4 \neq 2\), this is a counterexample, and therefore the claim is
  false.
\end{example}

\subsubsection{Proof by contradiction}

In a proof by contradiction, we assume that the claim is false and show that
this assumption leads to a contradiction.

\begin{example}
  Suppose that \(A \subseteq B\) and \(A \nsubseteq C\). We prove by
  contradiction that \(B \nsubseteq C\).

  Assume, to the contrary, that \(B \subseteq C\). Since
  \(A \nsubseteq C\), there exists an element \(x \in A\) such that
  \(x \notin C\). But \(A \subseteq B\) implies that \(x \in B\), and
  \(B \subseteq C\) then implies that \(x \in C\). This contradicts
  \(x \notin C\). Therefore \(B \nsubseteq C\).
\end{example}

\subsubsection{Relations}

\begin{definition}
  A \emph{relation} from a set \(A\) to a set \(B\) is a subset
  \[
    R \subseteq A \times B.
  \]
  If \((a,b) \in R\), we say that \(a\) is related to \(b\) and write
  \(aRb\). A relation on \(A\) is a relation from \(A\) to itself, that is, a
  subset of \(A \times A\).
\end{definition}

\begin{definition}
  Let \(R \subseteq A \times B\) be a relation. The \emph{inverse relation}
  of \(R\) is
  \[
    R^{-1} = \{(b,a) \in B \times A : (a,b) \in R\}.
  \]
  Equivalently, \(bR^{-1}a\) if and only if \(aRb\).
\end{definition}

\begin{definition}
  Let \(R\) be a relation on a set \(A\). The relation \(R\) is
  \begin{enumerate}
    \item \emph{reflexive} if \(xRx\) for every \(x \in A\);
    \item \emph{irreflexive} if \((x,x) \notin R\) for every \(x \in A\);
    \item \emph{symmetric} if, for all \(x,y \in A\),
      \(xRy\) implies \(yRx\);
    \item \emph{antisymmetric} if, for all \(x,y \in A\),
      \(xRy\) and \(yRx\) imply \(x=y\);
    \item \emph{asymmetric} if, for all \(x,y \in A\),
      \(xRy\) implies \((y,x) \notin R\);
    \item \emph{transitive} if, for all \(x,y,z \in A\),
      \(xRy\) and \(yRz\) imply \(xRz\); and
    \item \emph{total} if, for all \(x,y \in A\), either \(xRy\) or
      \(yRx\).
  \end{enumerate}
  A relation is an \emph{equivalence relation} if it is reflexive, symmetric,
  and transitive.
\end{definition}

\subsection{17.09.2026}

\subsubsection{Properties of inverse relations}

\begin{theorem}
  Let \(R,S \subseteq A \times B\) be relations. Then
  \begin{enumerate}[label=(\alph*)]
    \item \((R^{-1})^{-1} = R\);
    \item \((R \cup S)^{-1} = R^{-1} \cup S^{-1}\);
    \item \((R \cap S)^{-1} = R^{-1} \cap S^{-1}\);
    \item \((R \setminus S)^{-1} = R^{-1} \setminus S^{-1}\).
  \end{enumerate}
\end{theorem}

\begin{proof}
  For (a), let \((a,b) \in A \times B\). By the definition of the inverse
  relation,
  \[
    (a,b) \in (R^{-1})^{-1}
    \iff (b,a) \in R^{-1}
    \iff (a,b) \in R.
  \]
  For the remaining identities, let \((b,a) \in B \times A\). Then
  \begin{align*}
    (b,a) \in (R \cup S)^{-1}
    &\iff (a,b) \in R \cup S \\
    &\iff (a,b) \in R \lor (a,b) \in S \\
    &\iff (b,a) \in R^{-1} \lor (b,a) \in S^{-1} \\
    &\iff (b,a) \in R^{-1} \cup S^{-1},
  \end{align*}
  which proves (b). Similarly,
  \begin{align*}
    (b,a) \in (R \cap S)^{-1}
    &\iff (a,b) \in R \land (a,b) \in S \\
    &\iff (b,a) \in R^{-1} \land (b,a) \in S^{-1} \\
    &\iff (b,a) \in R^{-1} \cap S^{-1},
  \end{align*}
  proving (c). Finally,
  \begin{align*}
    (b,a) \in (R \setminus S)^{-1}
    &\iff (a,b) \in R \land (a,b) \notin S \\
    &\iff (b,a) \in R^{-1} \land (b,a) \notin S^{-1} \\
    &\iff (b,a) \in R^{-1} \setminus S^{-1},
  \end{align*}
  which proves (d). In each case, the two sets have exactly the same elements.
\end{proof}

\begin{remark}
  In (d), the handwritten notes have an extra inverse on \(S\) inside the
  parentheses. The corrected identity is stated above. If the expression
  \((R \setminus S^{-1})^{-1}\) is used literally, its value is
  \(R^{-1} \setminus S\).
\end{remark}

\subsubsection{Composition of relations}

\begin{definition}
  Let \(R \subseteq X \times Y\) and \(S \subseteq Y \times Z\).
  Their \emph{composite relation} is the relation from \(X\) to \(Z\)
  defined by
  \[
    S \circ R \coloneqq
    \set{(x,z) \in X \times Z :
    \exists y \in Y,\ (x,y) \in R \land (y,z) \in S}.
  \]
  Equivalently,
  \[
    x(S \circ R)z
    \iff \exists y \in Y : xRy \land ySz.
  \]
  Thus \(S \circ R\) means that we first follow \(R\), then \(S\).
\end{definition}

\begin{theorem}[Associativity of composition]
  Let \(R \subseteq X \times Y\), \(S \subseteq Y \times Z\), and
  \(T \subseteq Z \times U\). Then
  \[
    T \circ (S \circ R) = (T \circ S) \circ R.
  \]
\end{theorem}

\begin{proof}
  Let \((x,u) \in X \times U\). If
  \((x,u) \in T \circ (S \circ R)\), there is a \(z \in Z\) such that
  \(x(S \circ R)z\) and \(zTu\). The first relation supplies a
  \(y \in Y\) with \(xRy\) and \(ySz\). Since \(ySz\) and \(zTu\),
  we have \(y(T \circ S)u\), and hence
  \((x,u) \in (T \circ S) \circ R\).

  Conversely, if \((x,u) \in (T \circ S) \circ R\), there is a
  \(y \in Y\) such that \(xRy\) and \(y(T \circ S)u\). Choose
  \(z \in Z\) with \(ySz\) and \(zTu\). Then \(x(S \circ R)z\),
  so \((x,u) \in T \circ (S \circ R)\). Both inclusions hold, proving
  the equality.
\end{proof}

\subsubsection{Equivalence relations and equivalence classes}

Recall that a relation \(\sim\) on a set \(X\) is an
\emph{equivalence relation} if it is reflexive, symmetric, and transitive.

\begin{problem}[Lecture exercise]
  Let \(A = \set{2,3,4,5,6}\), and define a relation \(R\) on \(A\) by
  \[
    xRy \iff x \text{ and } y \text{ have a common prime factor}.
  \]
  Is \(R\) an equivalence relation?
\end{problem}

\begin{solution}
  The relation is reflexive: each \(x \in A\) has a prime factor, and that
  prime divides both occurrences of \(x\). More explicitly, \(2\) divides
  \(2,4,6\), \(3\) divides \(3\), and \(5\) divides \(5\).

  It is also symmetric. If \(xRy\), some prime \(p\) divides both \(x\)
  and \(y\). The same prime divides both \(y\) and \(x\), so \(yRx\).

  However, it is not transitive. We have \(4R6\), since \(2\) divides
  both \(4\) and \(6\), and \(6R3\), since \(3\) divides both \(6\)
  and \(3\). But \(4\) and \(3\) have no common prime factor: their
  only prime factors are \(2\) and \(3\), respectively. Thus
  \((4,3) \notin R\). Consequently, \(R\) is not an equivalence relation.
\end{solution}

\begin{definition}
  Let \(\sim\) be an equivalence relation on a set \(X\), and let
  \(a \in X\). The \emph{equivalence class} of \(a\) is
  \[
    [a]_{\sim} \coloneqq \set{x \in X : x \sim a}.
  \]
  The element \(a\) is called a \emph{representative} of this class.
  The set of all equivalence classes is the \emph{quotient set}
  \[
    X/{\sim} \coloneqq \set{[a]_{\sim} : a \in X}.
  \]
\end{definition}

The following consequence explains why equivalence classes divide a set into
separate parts; it supplements the definitions in the lecture.

\begin{proposition}
  Let \(\sim\) be an equivalence relation on \(X\). For every \(a,b \in X\),
  \begin{enumerate}[label=(\alph*)]
    \item \(a \in [a]_{\sim}\);
    \item \(a \sim b\) if and only if \([a]_{\sim} = [b]_{\sim}\);
    \item the classes \([a]_{\sim}\) and \([b]_{\sim}\) are either equal
      or disjoint.
  \end{enumerate}
  Therefore the equivalence classes form a \emph{partition} of \(X\):
  they are nonempty, their union is \(X\), and distinct classes are disjoint.
\end{proposition}

\begin{proof}
  Reflexivity gives \(a \sim a\), proving (a).

  For (b), suppose \(a \sim b\). If \(x \in [a]_{\sim}\), then
  \(x \sim a\), so transitivity gives \(x \sim b\). Hence
  \([a]_{\sim} \subseteq [b]_{\sim}\). Symmetry gives \(b \sim a\),
  and the same argument gives the reverse inclusion. Conversely, if the
  classes are equal, then (a) gives \(a \in [a]_{\sim} = [b]_{\sim}\),
  which means \(a \sim b\).

  For (c), if the classes intersect, choose an element \(c\) in their
  intersection. Then \(c \sim a\) and \(c \sim b\). By symmetry,
  \(a \sim c\), so transitivity gives \(a \sim b\). Part (b) now gives
  equality of the classes. Thus distinct classes must be disjoint.
  Finally, every class is a subset of \(X\), and (a) shows that each
  element of \(X\) belongs to a nonempty class. Their union is therefore
  all of \(X\), proving the partition claim.
\end{proof}

\subsubsection{Functions}

\begin{definition}
  Let \(X\) and \(Y\) be nonempty sets. A \emph{function} (or
  \emph{mapping}) \(f\) from \(X\) to \(Y\) assigns to each element
  \(x \in X\) exactly one element \(y \in Y\). We write
  \[
    f\colon X \longrightarrow Y,
    \qquad y = f(x).
  \]
  The set \(X\) is the \emph{domain}, and \(Y\) is the \emph{codomain}.
\end{definition}

In terms of relations, the graph of \(f\) is
\[
  G_f = \set{(x,f(x)) : x \in X} \subseteq X \times Y.
\]
A relation \(F \subseteq X \times Y\) is the graph of a function precisely
when, for every \(x \in X\), there exists exactly one \(y \in Y\) such
that \((x,y) \in F\). Existence ensures that the function is defined at
every input, and uniqueness ensures that its value is unambiguous.

\end{document}
